% Generated draft: requires mathematical and source review.
\documentclass[11pt]{article}
\usepackage[margin=25mm]{geometry}
\usepackage{fontspec}
\setmainfont{Latin Modern Roman}
\usepackage{amsmath,amssymb,mathtools,unicode-math}
\setmathfont{Latin Modern Math}
\usepackage{enumitem,needspace,xurl,hyperref}
\setlength{\parindent}{0pt}
\setlength{\parskip}{5pt}
\setlength{\emergencystretch}{3em}
\newcommand{\N}{\mathbb{N}}
\newcommand{\Z}{\mathbb{Z}}
\newcommand{\Q}{\mathbb{Q}}
\newcommand{\R}{\mathbb{R}}
\newcommand{\C}{\mathbb{C}}
\newcommand{\F}{\mathbb{F}}
\newcommand{\A}{\mathbb{A}}
\newcommand{\PP}{\mathbb{P}}
\newcommand{\E}{\mathbb{E}}
\newcommand{\eps}{\varepsilon}
\newcommand{\dd}{\,\mathrm{d}}
\begin{document}
\section*{UnsolvedMath ID: 30004837}
\textbf{Source problem code:} catalogue entry
\subsection{Definitions and mathematical statement}
The terms \emph{symmetric Macdonald polynomials} and \emph{their quasi-polynomial generalizations} refer to the objects intended in the cited Oberwolfach discussion. The supplied extract does not define the latter family or specify its root data, parameters, coefficient field, or parameter restrictions, so these choices cannot be fixed from this record alone.

\textbf{Research question.} Do the quasi-polynomial generalizations of symmetric Macdonald polynomials discussed in the cited source admit natural quantum-group interpretations?

Here an \emph{interpretation} would require a mathematical construction relating the intended quasi-polynomial family to quantum-group representation theory. The record specifies neither the quantum groups nor the required relationship, and it supplies no criterion for \emph{natural}. Consequently, this is a source-faithful research question rather than a fully specified existence proposition. In particular, the record does not determine whether the question concerns every member of a family, all parameter choices, or only particular cases. A precise formulation requires identifying the intended quasi-polynomial construction and specifying the representation-theoretic relationship and its hypotheses.
\subsection{Short English statement}
The question asks whether the particular quasi-polynomial extensions of symmetric Macdonald polynomials intended by the source have a natural explanation through quantum groups. The family, admissible parameters, and meaning of a natural interpretation require clarification before this becomes a precise mathematical proposition.
\subsection{Sources}
\begin{itemize}
\item Supplied local UnsolvedMath registry record: permanent ID \(30004837\), title \emph{Quantum-Group Interpretations of Quasi-Macdonald Polynomials}, source problem code \texttt{catalogue entry}, in the Oberwolfach Reports Open Problems collection. No individual proposer is supplied; the registry gives proposed year \(2021\).
\item \emph{Quantum Groups -- Algebra, Analysis and Category Theory} (2022), Oberwolfach Reports, pp.\ 2397--2458, DOI \texttt{10.4171/owr/2021/44}: \url{https://doi.org/10.4171/owr/2021/44}. The supplied original extract carries page number \(2437\). Oberwolfach Reports are described in the record as documenting mathematical workshops and their open-problem discussions. The citation year \(2022\) and the registry's proposed year \(2021\) are retained as supplied.
\item Supplied literature source: \emph{Quasi-polynomial DAHA representations}, DOI \texttt{10.1017/fms.2025.33}: \url{https://doi.org/10.1017/fms.2025.33}. The record supplies no authors or detailed construction from this reference.
\item Literature assessment dated exactly \texttt{2026-08-22}: the supplied curation labels the problem open and reports work on representations of double affine Hecke algebras (DAHA) and quasi-polynomial analogues of symmetric Macdonald polynomials, while describing natural quantum-group interpretations of these particular generalizations as still open. This is supplied dated curation, not an independently verified literature assessment; the record expressly requests independent verification of time-sensitive status and cited claims.
\item The registry labels the statement \texttt{corrected\_verified} and reports that the linked report was inspected and its explicit open question isolated in substance. This is supplied verification provenance; no source inspection or mathematical verification is asserted by this draft.
\end{itemize}
\end{document}
